\lysh{13090}{4}

\begin{alist}
\item Αφού το σημείο $M(x,y)$ κινείται στη γραφική παράσταση της συνάρτησης $f$ ισχύει ότι
$y = f(x) = \dfrac{16}{x} > 0$. Οι συντεταγμένες του σημείου Β είναι $B(x, 0)$ και του σημείου $A\left(0, \dfrac{16}{x}\right)$.
Το ορθογώνιο ΟΑΜΒ έχει εμβαδόν $(OAMB) = (OA) \cdot (OB) = \dfrac{16}{x} \cdot x = 16$ τετραγωνικές μονάδες και περίμετρο $\Pi(x) = 2(OA) + 2(OB) = 2 \cdot \dfrac{16}{x} + 2 \cdot x = 2x + \dfrac{32}{x}, x > 0$.

\item Αναζητούμε τις τιμές του $x > 0$ για τις οποίες $\Pi(x) = 20$, δηλαδή $2x + \dfrac{32}{x} = 20$ και ισοδύναμα $2x^2 + 32 = 20x$, δηλαδή $2x^2 - 20x + 32 = 0$ και τελικά $x^2 - 10x + 16 = 0$.
Η εξίσωση αυτή έχει ρίζες τις $x = 2$ ή $x = 8$ αφού $S = 2 + 8 = 10$ και $P = 2 \cdot 8 = 16$.
Για $x = 2$ είναι $y = \dfrac{16}{2} = 8$ οπότε $M_1(2, 8)$.
Για $x = 8$ είναι $y = \dfrac{16}{8} = 2$ οπότε $M_2(8, 2)$.

\item 
\begin{rlist}
\item Το $OAM'B$ είναι τετράγωνο αν και μόνο αν $(OA) = (OB)$ δηλαδή ισοδύναμα που σημαίνει $\dfrac{16}{x} = x$
οπότε $x^2 = 16$ και επειδή $x > 0$ έχουμε ότι $x = 4$.
\item Θα δείξουμε ότι $\Pi(x) \ge \Pi(4)$ για κάθε $x > 0$, δηλαδή ισοδύναμα
$2x + \dfrac{32}{x} \ge 16$ και επειδή $x > 0$ έχουμε ισοδύναμα $2x^2 + 32 \ge 16x$ δηλαδή $2x^2 + 32 - 16x \ge 0$
δηλαδή $x^2 + 16 - 8x \ge 0$ και τελικά $(x - 4)^2 \ge 0$ που ισχύει για κάθε $x > 0$, με την ισότητα να
ισχύει μόνο για $x = 4$, που σημαίνει ότι από όλα τα ορθογώνια $OAMB$ τη μικρότερη
περίμετρο την έχει το τετράγωνο $OAM'B$.
\end{rlist}
\end{alist}
